% !TeX program = xelatex
% Compile from this directory: latexmk -xelatex template.tex
\documentclass[12pt]{article}

% Page size and margins follow the supplied Word reference (US Letter).
\usepackage[letterpaper,left=1.25in,right=1.25in,top=1in,bottom=1in]{geometry}
\usepackage{fontspec}
\IfFontExistsTF{Times New Roman}
  {\setmainfont{Times New Roman}}
  {\setmainfont{TeX Gyre Termes}}
% Portable Chinese support, for example when filling in names.
\usepackage{xeCJK}
\setCJKmainfont{FandolSong-Regular.otf}[BoldFont=FandolSong-Bold.otf]
\usepackage{amsmath}
\usepackage{graphicx}
\usepackage{booktabs,tabularx}
\usepackage{caption}
\usepackage{placeins}
\usepackage{xcolor}
\usepackage{titlesec}
\usepackage[hidelinks]{hyperref}
\graphicspath{{fig/}}
\setlength{\parindent}{2em}
\setlength{\parskip}{0.3em}
\titleformat{\section}{\fontsize{14}{17}\selectfont\bfseries}{\thesection.}{0.6em}{}
\titleformat{\subsection}{\normalsize\bfseries}{\thesubsection}{0.6em}{}
\captionsetup{font=small,labelfont=bf,labelsep=period}
\definecolor{courseblue}{HTML}{0070C0}

% ---------- Edit these cover fields ----------
\newcommand{\University}{Southern University of Science and Technology}
\newcommand{\Department}{Department of Materials Science and Engineering}
\newcommand{\Course}{Experiments for Fundamental of Materials Science and Engineering MSE002}
\newcommand{\LabNumber}{[Number]}
\newcommand{\ExperimentTitle}{[Experiment title]}
\newcommand{\StudentName}{[Your name]}
\newcommand{\StudentID}{[Your student ID]}
\newcommand{\TeamMembers}{[Team member names]}
\newcommand{\ExperimentDate}{[YYYY/MM/DD]}
\newcommand{\SubmissionDate}{[YYYY/MM/DD]}

\hypersetup{pdftitle={Lab Report Template},pdfauthor={},pdfsubject={Reusable laboratory report template}}

\begin{document}

% ==============================================================================
%                           封面部分开始 / COVER START
% ==============================================================================
% Editable cover: no external PDF or personal information is required.
\hypersetup{pageanchor=false}
\begin{titlepage}
  \thispagestyle{empty}
  {\noindent\fontsize{15}{18}\selectfont\bfseries
   \University\par
   \noindent\Department\par}
  \vspace{1.7cm}
  \begin{center}
    {\fontsize{18}{22}\selectfont\bfseries\color{courseblue}\Course\par}
    \vspace{1.7cm}
    {\fontsize{28}{34}\selectfont\bfseries Lab Report\par}
    \vspace{0.6cm}
    {\fontsize{22}{27}\selectfont\bfseries Lab \LabNumber: \ExperimentTitle\par}
  \end{center}
  \vspace{1.8cm}
  \begin{flushright}
    \fontsize{15}{21}\selectfont\bfseries
    Name: \StudentName\par
    Student ID: \StudentID\par
    Team members: \TeamMembers\par
    Date of Experiment: \ExperimentDate\par
    Date of submission: \SubmissionDate\par
  \end{flushright}
  \vfill
\end{titlepage}
\setcounter{page}{1}
\hypersetup{pageanchor=true}

% ==============================================================================
%                           封面部分结束 / COVER END
% ==============================================================================

\section{Objectives}
State the purpose of the experiment and identify the quantities, relationships,
or techniques to be investigated. Replace the following items with specific,
measurable objectives.
\begin{enumerate}
  \item Describe the scientific question and the property to be measured.
  \item Identify the experimental method and the skill to be developed.
  \item State the comparison or hypothesis to be evaluated.
\end{enumerate}

\section{Principles}
Explain the theory needed to interpret the experiment in your own words.
Define the quantities and assumptions, and cite the sources used~\cite{reference}.
Include only equations relevant to the analysis. For example, the sample mean is
\begin{equation}
  \overline{x}=\frac{1}{n}\sum_{i=1}^{n}x_i,
  \label{eq:mean}
\end{equation}
where $x_i$ is the $i$th measurement and $n$ is the number of measurements.
Replace or extend Eq.~\eqref{eq:mean} with the equations used in your experiment.

\begin{figure}[htbp]
  \centering
  \includegraphics[width=0.62\linewidth]{figA.jpg}
  \caption{Replace Fig A with a diagram of the experimental principle.
  Define the components and symbols in the caption.}
  \label{fig:principle}
\end{figure}

\section{Experimental (Materials and Apparatus)}
\subsection{Materials}
List the materials, specimen identifiers, compositions, dimensions, and preparation
or storage conditions needed to reproduce the experiment. Specify grades,
purities, and suppliers where relevant.

\subsection{Apparatus}
Identify each instrument, its model, measurement range, resolution, and calibration
status where relevant. Explain the arrangement shown in Fig.~\ref{fig:apparatus}.

\begin{figure}[htbp]
  \centering
  \includegraphics[width=0.62\linewidth]{figB.jpg}
  \caption{Replace Fig B with a photograph or schematic of the apparatus.
  Identify the instrument and the measurement arrangement.}
  \label{fig:apparatus}
\end{figure}

\section{Experimental procedure}
Describe what was actually done in chronological order. Include operating
conditions, calibration, repeated measurements, and any deviations from the
laboratory instructions.
\begin{enumerate}
  \item Describe specimen preparation and the initial instrument checks.
  \item Specify the controlled variables, settings, and measurement sequence.
  \item Explain how observations were recorded and how many repeats were made.
  \item Describe shutdown, cleaning, and disposal as applicable.
\end{enumerate}

\section{Results and discussion}
\textbf{This is a key section of the report.} Present the measurements and explain
what they show. Use descriptive captions and refer to every figure and table in
the text. Replace all placeholder images and entries with your own results.

\subsection{Results and calculations}
Organize measured and calculated quantities with units and suitable precision.
Table~\ref{tab:results} illustrates a results table; bracketed entries are fields
to complete, not measurements. Show the calculation method and report uncertainty
or variability where appropriate.
\begin{table}[htbp]
  \centering
  \caption{Replace this caption with the measured quantity and test conditions.}
  \label{tab:results}
  \begin{tabularx}{\linewidth}{@{}lXXX@{}}
    \toprule
    Sample & Condition [unit] & Mean [unit] & SD [unit] \\
    \midrule
    {[Sample 1]} & [Condition] & [Mean] & [SD] \\
    {[Sample 2]} & [Condition] & [Mean] & [SD] \\
    \bottomrule
  \end{tabularx}
\end{table}

\begin{figure}[htbp]
  \centering
  \includegraphics[width=0.62\linewidth]{figC.jpg}
  \caption{Replace Fig C with a results plot. Label the axes with quantities
  and units; define symbols, sample size, and error bars where used.}
  \label{fig:results}
\end{figure}

\subsection{Interpretation and sources of error}
Explain the trends in Fig.~\ref{fig:results} using the principles above. Compare
with expected behavior or cited literature and discuss disagreements. Identify
specific sources of error, their likely effects, and practical improvements.
Use Fig.~\ref{fig:comparison} for an additional comparison or observation if needed.

\begin{figure}[htbp]
  \centering
  \includegraphics[width=0.62\linewidth]{figD.jpg}
  \caption{Replace Fig D with a comparison plot, specimen image, or other
  observation that supports the discussion. Explain the relevant features.}
  \label{fig:comparison}
\end{figure}
\FloatBarrier

\section{Conclusion}
\textbf{This is a key section of the report.} Summarize the main findings and
state whether the objectives were achieved. Include the most relevant numerical
results with units and uncertainty when available. State the main limitation;
keep the conclusions supported by the results already presented.

\section{Questions}
\textbf{This is a key section of the report.} Reproduce each assigned question
and answer it directly. Include reasoning, calculations, or references as needed.
\begin{enumerate}
  \item \textbf{[Insert the assigned question.]}
  \par\noindent Answer: [Write your answer and supporting explanation here.]
  \item \textbf{[Insert the next assigned question.]}
  \par\noindent Answer: [Write your answer and supporting explanation here.]
\end{enumerate}

% Optional references: replace the placeholder below with a real source.
% Remove both this environment and its citation if no references are used.
\begin{thebibliography}{9}
  \bibitem{reference}
  [Author(s)]. [Title]. \textit{[Journal or publisher]}, [year],
  [volume(issue)], [pages]. [DOI or URL, if available].
\end{thebibliography}

% Keep attachments at the END, as requested by the Word reference.
\clearpage
\appendix
\section{Original datasheet}
Attach the original experimental datasheet at the end of the report. Replace
the table below with the actual records, or insert readable scans of all required
pages. Preserve sample identifiers, measurement units, test conditions, dates,
and signatures where required.

\begin{table}[htbp]
  \centering
  \caption{Replace this table with the original measurement records.}
  \begin{tabularx}{\linewidth}{@{}lXXX@{}}
    \toprule
    Record & Sample ID & Reading [unit] & Conditions / notes \\
    \midrule
    {[Record 1]} & [ID] & [Reading] & [Notes] \\
    {[Record 2]} & [ID] & [Reading] & [Notes] \\
    \bottomrule
  \end{tabularx}
\end{table}

\end{document}
